% ==============================================================================
% Präambel der Studienarbeit: Python, OCaml und F# in Stroetmanns Vorlesung
% ==============================================================================
\usepackage[utf8]{inputenc}
\usepackage[T1]{fontenc}
\usepackage[english,ngerman]{babel}

% Mathematik (amsthm muss vor newtxmath geladen werden)
\usepackage{amsmath,amsthm}
\usepackage{mathtools}

% Schrift: Times-artige Grundschrift, serifenlose Überschriften (KOMA-Standard)
\usepackage{newtxtext}
\usepackage{newtxmath}
\usepackage[scaled=0.88]{beramono}
\usepackage[htt]{hyphenat}
\usepackage{microtype}

% Satzspiegel
\usepackage[
  a4paper,
  top=2.5cm,
  bottom=2.5cm,
  left=3.0cm,
  right=2.5cm,
  headheight=15pt
]{geometry}

\usepackage{setspace}
\onehalfspacing

% Überschriften und Abstände
\addtokomafont{disposition}{\color{black}}
\addtokomafont{chapter}{\Huge}
\RedeclareSectionCommand[beforeskip=0pt,afterskip=1.5\baselineskip]{chapter}
\setkomafont{caption}{\small}
\setkomafont{captionlabel}{\small\bfseries}

% Tabellen
\usepackage{makecell}
\renewcommand\theadfont{\bfseries\small}
\renewcommand\theadalign{lc}
\usepackage{booktabs}
\usepackage{tabularx}
\usepackage{longtable}
\usepackage{multirow}
\usepackage{multicol}

% Listen
\usepackage{enumitem}
\setlist{itemsep=2pt,topsep=4pt}
\setlist[description]{style=nextline,leftmargin=1.2em,itemsep=6pt}

% Abbildungen
\usepackage{graphicx}
\usepackage{float}
\usepackage{caption}
\captionsetup{format=hang,labelsep=colon,justification=justified,singlelinecheck=true}
\usepackage{subcaption}
\usepackage{siunitx}
\sisetup{
  detect-all,
  group-separator={,},
  group-minimum-digits=4
}

% Farben: nur in Diagrammen, um die drei Sprachen zu unterscheiden
\usepackage[dvipsnames,table]{xcolor}
\definecolor{NavyBlue}{RGB}{40, 40, 40}
\definecolor{PythonBlue}{RGB}{48, 105, 152}
\definecolor{OCamlOrange}{RGB}{217, 95, 2}
\definecolor{FSharpPurple}{RGB}{117, 112, 179}
\definecolor{CodeBg}{RGB}{255, 255, 255}
\definecolor{CommentGray}{RGB}{90, 90, 90}

% Diagramme
\usepackage{tikz}
\usepackage{pgfplots}
\pgfplotsset{compat=1.18}
\usetikzlibrary{arrows.meta,positioning,shapes.geometric,calc,backgrounds,fit}

% Quelltexte
\usepackage{listings}

\lstdefinelanguage{FSharp}{
  morekeywords={
    let, rec, mutable, module, open, type, of, match, with, when, function,
    fun, if, then, else, elif, for, in, do, while, yield, return, try,
    finally, raise, failwith, true, false, null, and, or, not, as, begin,
    end, struct, val, member, interface, inherit, abstract, override, default,
    static, private, internal, public, inline, new
  },
  sensitive=true,
  morecomment=[l]{//},
  morecomment=[s]{(*}{*)},
  morestring=[b]"
}

% Schwarz-weiß: Schlüsselwörter fett, Kommentare kursiv, Linien oben und unten
\lstdefinestyle{academiccode}{
  basicstyle=\ttfamily\footnotesize\linespread{1.0}\selectfont,
  breakatwhitespace=false,
  breaklines=true,
  captionpos=b,
  commentstyle=\itshape\color{CommentGray},
  keywordstyle=\bfseries,
  stringstyle=\color{black},
  frame=lines,
  framesep=3pt,
  rulecolor=\color{black},
  keepspaces=true,
  numbers=left,
  numbersep=6pt,
  numberstyle=\tiny\color{CommentGray},
  showspaces=false,
  showstringspaces=false,
  showtabs=false,
  tabsize=2,
  xleftmargin=14pt,
  xrightmargin=0pt
}

\lstset{style=academiccode}

% Erlaubt zusätzlichen Wortabstand, wo langer Inline-Code sonst überstehen würde.
\setlength{\emergencystretch}{2em}

% Unicode-Zeichen der Notizbücher, die auch im Fließtext vorkommen
\DeclareUnicodeCharacter{22C5}{\ensuremath{\cdot}}
\DeclareUnicodeCharacter{2205}{\ensuremath{\emptyset}}
\DeclareUnicodeCharacter{1D700}{\ensuremath{\varepsilon}}
\DeclareUnicodeCharacter{1D6FF}{\ensuremath{\delta}}
\DeclareUnicodeCharacter{03A3}{\ensuremath{\Sigma}}

% Wörtlich aus den Notizbüchern übernommener Code: kleinere Schrift, damit drei
% Fassungen in eine Abbildung passen, und die Unicode-Zeichen der Notizbücher.
\lstdefinestyle{notebook}{
  style=academiccode,
  basicstyle=\ttfamily\fontsize{7.2}{8.4}\linespread{1.0}\selectfont,
  captionpos=t,
  aboveskip=2pt,
  belowskip=4pt,
  extendedchars=true,
  inputencoding=utf8,
  literate={𝜀}{{$\varepsilon$}}1 {⋅}{{$\cdot$}}1 {∅}{{$\emptyset$}}1
           {Σ}{{$\Sigma$}}1 {𝛿}{{$\delta$}}1
}

% Befund- und Beobachtungsumgebungen
\theoremstyle{definition}
\newtheorem{definition}{Definition}[chapter]
\newtheorem{observation}{Beobachtung}[chapter]
\newtheorem{finding}{Befund}[chapter]
\newtheorem{criterion}{Evaluierungskriterium}[chapter]

% Kopf- und Fußzeilen
\usepackage{scrlayer-scrpage}
\pagestyle{scrheadings}
\clearpairofpagestyles
\ihead{\headmark}
\ohead{\pagemark}
\automark[chapter]{chapter}
\KOMAoptions{headsepline=0.4pt}
\setkomafont{pageheadfoot}{\normalfont\small}

% Hyperlinks und Querverweise
\usepackage[
  hidelinks,
  bookmarksnumbered=true,
  pdftitle={Sprachparadigmen, Typsysteme und Ausfuehrungsdynamik: Python, OCaml und F-Sharp anhand der Vorlesung Formale Sprachen von Karl Stroetmann}
]{hyperref}
\usepackage[ngerman,nameinlink,noabbrev]{cleveref}
